\section{Kết luận và định hướng phát triển}
\label{chap:ket-luan}

\subsection{Tổng kết kết quả}
Trên official UNSW-NB15 TEST gồm 82\,332 dòng, Stacking-42 đạt Macro-F1 khoảng 0.863840 so với 0.863343 của RF-42, trong khi MCC của Stacking-42 khoảng 0.750497, thấp hơn rất nhẹ so với 0.750910 của RF-42. Chênh lệch Macro-F1 Stacking-42 trừ RF-42 là khoảng $+0.000497$; khoảng tin cậy 95\% paired bootstrap đã lưu gồm 0. Những bằng chứng này \emph{không xác nhận ưu thế tổng thể} của Stacking-42 trên Macro-F1/MCC, và cũng không chứng minh hai mô hình tương đương về mặt thống kê.

Stacking-42 giảm 142 dự đoán báo động giả (FP) nhưng tăng 118 trường hợp bỏ sót tấn công (FN) so với RF-42. Phân tích hậu nghiệm cho thấy xu hướng lỗi khác nhau theo từng attack family: Stacking-42 có thêm 133 FN ở Fuzzers nhưng giảm 43 FN ở Generic. Những con số đó thuộc mô hình \emph{binary} và nhãn thật của các nhóm, không phải khả năng dự đoán multiclass hay bằng chứng về nguyên nhân từ các feature riêng lẻ.

RF-17 có Macro-F1 thấp hơn RF-42 một lượng nhỏ, đổi lại các số đo lịch sử cho thấy thời gian train ngắn hơn và ước lượng dung lượng mô hình giảm. Tuy nhiên, đây là so sánh \emph{pipeline-level} vì RF-17 không chỉ giảm đặc trưng mà còn thay đổi siêu tham số Random Forest. Không thể quy trực tiếp mọi thay đổi cho riêng Feature Selection.

\subsection{Giải đáp các câu hỏi nghiên cứu}
\textbf{RQ1:} Dữ liệu official không cho thấy Stacking-42 vượt RF-42 rõ rệt về paired Macro-F1/MCC; việc bớt FP nhưng tăng FN là đánh đổi, còn ổn định repeated training chưa đo.
\textbf{RQ2:} RF-17 giảm chi phí lịch sử và thay đổi hiệu năng so với RF-42, nhưng vì cùng thay hyperparameters nên không kết luận hiệu ứng riêng của Feature Selection.
\textbf{RQ3:} Dữ liệu OOF cho thấy sai số bổ sung giữa base learners; việc Voting không tốt hơn RF cho thấy diversity không tự dẫn đến gain, và chưa chứng minh cơ chế nhân quả.
\textbf{RQ4:} Macro-F1 giảm ở novel-vector subset của cả ba pipeline nhưng thứ hạng giữ nguyên; kết quả không đại diện external validation.
\textbf{RQ5:} Mất cân bằng support và đánh đổi FPR/FNR có thể quan sát; do không có matched resampling/weighting study nên chưa xác định lợi ích riêng của cách xử lý mất cân bằng.
Các số liệu, bảng và mức giảm cụ thể nằm trong Mục~\ref{sec:revision-rq-all}; việc chưa thể xác nhận một giả thuyết \emph{không} được chuyển thành khẳng định thành công.

\subsection{Đóng góp trong phạm vi đã chứng minh}
Đóng góp có thể bảo vệ ở thời điểm hiện tại là một chuỗi thực nghiệm trên benchmark có nguồn và cấu hình đóng băng, có kiểm toán dữ liệu, đối chiếu mô hình học máy đơn và ensemble, phân tích sai số bổ sung, khoảng tin cậy cho chênh lệch theo cặp, đánh đổi FP/FN và chi phí. Kết quả không thuận lợi của Voting và sự thiếu ưu thế rõ ràng của Stacking là thành phần của kết luận khoa học, không bị loại vì không hỗ trợ một giả thuyết chiến thắng.

Báo cáo cũng ghi nhận các giới hạn của phép đo: sai lệch tiềm tàng do Feature Selection trước development CV, meta tuning chưa fully nested, vector TRAIN/TEST bị trùng, benchmark flow có thể phụ thuộc lẫn nhau, thiếu temporal/external validation và thiếu bằng chứng ổn định qua các lần tái huấn luyện độc lập. Việc công khai những giới hạn này góp phần tránh suy diễn vượt mức về khả năng tổng quát hoặc tính mới của thuật toán.

\subsection{Công việc còn thiếu và hướng nghiên cứu}
Với RQ1, cần một nghiên cứu mới có nhiều lần huấn luyện độc lập, cùng tài nguyên và môi trường đã khóa để kiểm tra tính ổn định của hiệu số. Với RQ2, cần thí nghiệm matched và nested feature selection, giữ các tham số còn lại cố định để tách tác động của chính Feature Selection. Với RQ3, có thể kiểm tra giả thuyết về cơ chế bổ sung sai số bằng phân tích độc lập thay vì chỉ dùng profile hậu nghiệm. Với RQ4, cần nguồn có thông tin thời gian/session hoặc một external benchmark đáng tin cậy để đánh giá ngoài official split. Với RQ5, có thể thiết kế đối chứng resampling/class-weight bên trong training folds, nhưng không được dựa vào official TEST hiện tại để chọn phương án thắng.

Mã nguồn xử lý, huấn luyện và tổng hợp chỉ số tạo điều kiện kiểm tra quy trình thực nghiệm và tính nhất quán của các kết quả. Môi trường huấn luyện lịch sử vẫn chưa đủ thông tin để xác nhận \emph{bitwise identity}. Tại checkpoint ngày 17/10/2026, ba frozen pipeline đã được tái dựng trên toàn 175\,341 dòng official TRAIN và lưu thành ba artifact \texttt{RECONSTRUCTED\_FULL\_TRAIN}; các phép xác minh trên laptop tác giả đạt 69/69 (SHA-256, schema, tải model và suy luận trên 16 dòng TRAIN). Đây là kết quả tái dựng có phạm vi, \emph{không} xác nhận model gốc, không tái đánh giá scientific TEST và không chứng minh hệ thống NIDS sẵn sàng triển khai. Các smoke model trên TRAIN nhỏ vẫn chỉ giữ vai trò thử kỹ thuật.

\subsection{Kết luận cuối}
Trong phạm vi benchmark và thí nghiệm đã khóa, kết quả chủ yếu nhấn mạnh nhu cầu so sánh Ensemble Learning một cách có kiểm soát, trình bày đầy đủ loại lỗi và độ bất định, thay vì chọn mô hình chỉ bởi một điểm Accuracy hoặc Macro-F1 nhỉnh hơn về mặt số học. Việc bổ sung Stacking đem lại một đánh đổi FP/FN cụ thể nhưng không xác lập ưu thế tổng thể trước RF-42. Các đóng góp có thể truy vết nằm ở cách đánh giá và diễn giải kết quả, không phải một cam kết superiority, novelty hoặc deployment chưa được kiểm chứng.
